\newpage
\section{Introducción}
\paragraph{}
En este trabajo de laboratorio se lleva a cabo el análisis de funcionamiento del un amplificador real para lo cual se tiene en cuenta las fuentes de error presentes en el mismo, causantes de que la salida final presente una desviación respecto a la respuesta ideal teórica.
\paragraph{}
El circuito sobre el cual se lleva a cabo el estudio es un sumador inversor. 
\paragraph{}
Los errores sobre los cuales se realiza el análisis son:
\begin{itemize}
    \item \textbf{Errores en DC.}
    \begin{enumerate}
	\item Error de tensión, $v_{os}$.
	\item Error de corriente, $I_{os}$.
	\item Error por $A_d<\infty$.
        \item Error por $RRMC<\infty$.
    \end{enumerate}
    \item \textbf{Errores en AC.}
    \begin{enumerate}
        \item Error vectorial.
        \item Ancho de banda a plena potencia.
        \item Ancho de banda para pqeuña señal.
    \end{enumerate}
\end{itemize}
\paragraph{}
Para cada circuito, se realizará un análisis teórico, simulaciones y mediciones experimentales. Finalmente, se van a comparar los datos obtenidos en cada etapa.

\section{Objetivos}
\begin{itemize}
	\item Aplicar el conocimiento teórico - práctico para analizar los errores acarreados en circuitos que hacen uso de amplificadores operacionales.
        \item Identificar las limitaciones respecto a la exactitud del respuesta ofrecida por el amplificador utilizado.
	\item Fortalecer el uso del simulador LtSpice e interpretar los resultados del mismo realizando las comparaciones correspondientes con los resultados teóricos obtenidos.
\end{itemize}
